\chapter{טעויות נפוצות וניהול סיכונים}

ניתוח של מאות כשלונות מצביע על דפוסים חוזרים. הסיבה המובילה לכישלון היא
היעדר צורך אמיתי בשוק, ואחריה אזילת המזומן וצוות לא מתאים~\cite{cbinsights2021}.
איור~\ref{fig:funnel} מציג את שרשרת הסיכונים העיקרית.

\begin{figure}[h]
  \centering
  % English labels + explicit coordinates inside an LTR (\entikz) context, so
  % babel bidi keeps the labels in their boxes. Hebrew meaning is in the caption.
  \entikz{%
  \begin{tikzpicture}[
      box/.style={draw=brandred, rounded corners, minimum width=2.6cm, minimum height=1.1cm,
                  align=center, font=\small, fill=brandred!8},
      arr/.style={-{Stealth[length=3mm]}, very thick, brandred}]
    \node[box] (need) at (0,0)   {No market\\need};
    \node[box] (cash) at (3.3,0) {Ran out\\of cash};
    \node[box] (team) at (6.6,0) {Wrong\\team};
    \node[box] (comp) at (9.9,0) {Out-\\competed};
    \draw[arr] (need) -- (cash);
    \draw[arr] (cash) -- (team);
    \draw[arr] (team) -- (comp);
  \end{tikzpicture}}
  \caption{שרשרת הסיכונים השכיחה לכישלון סטארט-אפ — היעדר צורך בשוק, אזילת מזומן,
  צוות לא מתאים ותבוסה בתחרות (לפי \en{CB Insights}; שורטט ב-\en{TikZ}).}
  \label{fig:funnel}
\end{figure}

\section{אנטי-דפוסים שכדאי להכיר}

\begin{itemize}
  \item \textbf{בנייה מוקדמת מדי:} כתיבת קוד לפני אימות הביקוש — בזבוז המשאב
        היקר ביותר, זמן.
  \item \textbf{הרחבה מוקדמת \term{(premature scaling)}:} השקעה בשיווק וגיוס
        לפני שהושגה התאמת מוצר-שוק.
  \item \textbf{מדדי יוהרה:} מעקב אחר מספרים מרשימים שאינם מנחים החלטה.
  \item \textbf{התעלמות מ-\en{runway}:} אי-ניהול קצב השריפה עד שמתגלה מאוחר
        מדי שאין מספיק מזומן לסבב הבא.
\end{itemize}

\section{ניהול סיכונים כתהליך מתמשך}

ניהול סיכונים אינו אירוע חד-פעמי אלא בקרה שוטפת: לזהות את ההשערה המסוכנת ביותר
בכל רגע, לבדוק אותה בזול, ולעדכן את התוכנית בהתאם. זוהי בדיוק לולאת הלמידה
מפרק~\ref{fig:bml}, מיושמת על קיום החברה עצמה.
